% An invented manuscript for trying Sextant. Every author, work, figure and
% DOI in it is made up; the DOIs use Crossref's 10.5555 test prefix.
\documentclass[11pt]{article}
\usepackage[T1]{fontenc}
\usepackage[margin=2.6cm]{geometry}
\usepackage{lmodern}
\usepackage{xcolor}
\usepackage[colorlinks=true,citecolor=blue!60!black,linkcolor=black,urlcolor=black]{hyperref}

\title{Raised Beds, Mulch and Compost in Shared Urban Gardens:\\ A Field Note}
\author{Ada Example \and Ben Placeholder\\[2pt]
  \small Department of Imaginary Studies, Nowhere University}
\date{}

\begin{document}
\maketitle

\begin{abstract}
Community gardens on reclaimed city lots start with dry, compacted soil and
little money to fix it. We compare three inexpensive interventions --
raised beds, straw mulch and compost -- using results already reported in the
literature, and ask which of them a volunteer group should try first.
\end{abstract}

\section{Introduction}
Most shared gardens are built on land that was never meant to grow anything.
Builders' rubble, years of foot traffic and a thin layer of imported topsoil
leave plots that dry out within days of rain and crust over in summer.
Volunteers rarely have the budget for irrigation, so the question is which
low-cost change does the most for soil water.

\section{What the literature reports}
Raised beds are the most visible intervention. Over three growing seasons,
raised beds retained 23\% more soil moisture than in-ground plots
\cite{example2021beds}, largely because their loose fill holds water that a
compacted plot sheds as run-off.

Mulch addresses the surface instead. A straw mulch layer 8 cm deep cut weed
emergence by 54\% while slowing evaporation \cite{placeholder2019mulch}, and
it costs little more than the straw itself.

Compost works on the soil structure underneath. With compost amendments,
compacted plots absorbed 31 mm of simulated rainfall in the first hour
\cite{testwell2022compost}, against a fraction of that on untreated ground.

\section{Discussion}
The three interventions are complementary rather than competing. Compost
improves infiltration, mulch keeps that water from evaporating, and raised
beds make both easier to maintain on a site that cannot be dug. For a group
starting from bare ground, compost is the cheapest first step; mulch follows
naturally once there is something to protect; raised beds are worth their cost
where the ground is contaminated or too compacted to work at all.

\section{Conclusion}
None of these figures is new; they come from the works cited. The point of
this note is the comparison, and every number in it can be traced to the
passage that reports it.

\bibliographystyle{unsrt}
\bibliography{refs}
\end{document}
